\section{Introduction}
\label{sec:intro}
% src: TEST_RESULTS.md (gate), TEST_RESULTS_HUMAN_SYM.md, posthoc_2026-11/{S2a_packs,S6_baselines,S3_agentic,S5_stats,S7_analysis} (see Sections 4-5 for the numbers)
% src (REVIEW2 N1, 2026-11-11): 15 to 18 and 11 to 15 = funnel of Table funnel (S10_retrieval sec. 2.4; S7_analysis sec. 1, 5); 5 to 10 at 6,000 = S8_factorial (+9..+11 saved answers) and S14_bilingual sec. 2.1 (+4..+6, +4..+5 redraws); 15 at 3,000 = S8 (+15..+16, one run);
%      15 to 18 fewer answers with translated questions = S15_crossling sec. 2 (HR 69..71 -> 53..54); SD about two, up to six = S2b_variance Table variance.
% Reframed 2026-11-11 to the diagnostic framing of work/review_fable/FRAMING.md (F1): where answers are lost; the guide is the closing synthesis. No bullet lists; numbers kept to the few that carry a claim.

Anyone who builds a retrieval-augmented generation (RAG) system has to make a series of choices. The search can be lexical (BM25), dense (vector similarity) or a hybrid of the two. The candidates can then be reranked by a cross-encoder or selected by a language model. The system must decide how much text the reader model sees, which we call the \emph{evidence budget}, meaning the number of characters of retrieved text in the prompt. It can search once or iteratively, or skip search and read long documents. There are also smaller choices of chunk size, query rewriting and embedding model. Method papers usually vary one of these at a time on English benchmarks. Component surveys and cost comparisons exist \citep{wang2024bestpractices,li2024selfroute,ferrazzi2026agentic}, but a practitioner who works in Persian, or with a tight token budget, still has little to go on.

This paper asks where a RAG pipeline loses its answers. Our finding is that they are lost mostly before the reader, in two stages. For 15 to 18 of the 90 test questions the 40 top-ranked candidates held no complete evidence bundle, that is, no set of passages that together support the answer. For another 11 to 15, a complete bundle was among the candidates and was dropped when the evidence pack, the text shown to the reader, was built. The reader itself is rarely the cause. A reranking or selection step, a larger budget and iterative search each recover part of the loss under different conditions, which we describe. The closest study of tight budgets found that what survives packing into the context predicts accuracy better than recall does \citep{bala2026recall}. Our results agree and extend that view to natural documents and whole pipelines.

We examine the question in Persian and in English. Persian stresses lexical retrieval, since Arabic-script variants, zero-width non-joiners and mixed digit systems all affect matching. We built a benchmark with company reports, books and statutes in each language. Its questions come in three tiers. A tier-1 (T1) question asks for a single fact. A tier-2 (T2) question needs evidence combined within one document, and a tier-3 (T3) question needs evidence from at least two documents. The two languages use different documents, so the Persian and English halves are two runs of one design and not a controlled contrast between languages (Section~\ref{sec:benchmark} says what the design can and cannot show).

The comparisons start with a frozen one-shot test of four whole pipelines at two evidence budgets, 6,000 and 15,500 characters (Figure~\ref{fig:pipelines}). The baseline, \hybrid, combines BM25 and dense retrieval, and \hr{} adds a cross-encoder reranker. Our method, \lexb{}, fuses three retrieval signals and then makes one LLM call that names bundles of passages that together support the answer. The fourth pipeline, \setr, is an adaptation of the published SetR selector \citep{lee2025setr}, which introduced requirement-driven set selection; \lexb{} is a zero-shot, budget-aware relative of it. This frozen test is the only confirmatory part of the paper. \lexb{} did not reach the margin we had fixed in advance over both controls, and no pipeline was clearly better than the cross-encoder reranker.

That outcome left the diagnosis open, so we used the same test set for a second study and label every result in it as post hoc. We first vary the retrieval components one at a time, namely the retrieval channel, query rewriting, chunk size and the embedding model. Hybrid retrieval is clearly better than BM25 alone, and none of the other changes beat the default. We then ask where the answers are lost. With only the gold evidence the reader answers 83 of 90 questions, against 26 with no evidence at all, so most misses come from evidence that did not reach it. At 6,000 characters, which we call a tight budget, removing Lex-RAG's selection call costs 5 to 10 answers, and a cross-encoder reaches a similar score with no LLM tokens. A bigger pack lets plain hybrid retrieval catch up, but at several times the evidence, and the reranked pipeline gains less from it.

Searching harder has the best point estimate. An iterative agent that issues several queries scores highest at about four times the tokens, but its lead over the reranker is within the noise of the test. A separate subsection returns to the two languages. Persian costs more tokens per character, and when questions are asked in the other language BM25 fails and even the best pipeline answers 15 to 18 fewer questions. Last, we ask how far any of these rankings can be trusted. The reader's answers vary with a standard deviation of about two questions per cell between identical runs, and single cells moved by up to six. Two LLM judges order the middle pipelines differently, and a human review of only some systems' answers once produced a significant result that disappeared when every system was reviewed. We therefore give each recommendation with a confidence label.

The paper contributes a Persian-English benchmark with annotated evidence bundles, a frozen confirmatory test with its null result, a post hoc account of where answers are lost and which steps recover them, and measurements of token cost, memory and ranking stability. Section~\ref{sec:methods} describes the benchmark, the pipelines, the readers and judge, and the protocol. Section~\ref{sec:results} reports the results. Its first part, the frozen test (Section~\ref{sec:frozen}), is the only confirmatory one. The post hoc parts cover retrieval components (Section~\ref{sec:components}), where answers are lost (Section~\ref{sec:lost}), how much evidence to show (Section~\ref{sec:budget}) and whether more search helps (Section~\ref{sec:search}). They are followed by the Persian and English halves (Section~\ref{sec:language}) and by recommendations for choosing a setup (Section~\ref{sec:practice}). Section~\ref{sec:stability} checks how far the rankings can be trusted, and Section~\ref{sec:conclusion} concludes with what would settle the question. We will release the benchmark with gold evidence bundles, per-question verdicts and the full record of runs, including those that failed.
